\section{Conclusi\'on}

El desarrollo de la librer\'ia \texttt{orange-jumpsuit}
ha permitido validar que es posible implementar un motor de b\'usqueda
vectorial funcional y competitivo bajo una filosof\'ia de desarrollo de bajo
nivel y m\'inimas dependencias en relativamente poco tiempo. 
A lo largo de este trabajo, se ha logrado transformar los conceptos te\'oricos
de la cuantificaci\'on de productos en una herramienta pr\'actica que equilibra
la complejidad algor\'itmica con la eficiencia en sistemas modernos.

Los resultados obtenidos en las pruebas comparativas contra FAISS demuestran
que la implementación alcanza niveles de precisi\'on (Recall) similares
en conjuntos de datos est\'andar como SiftSmall, logrando un 83\% de
efectividad en la recuperación de los 10 vecinos más cercanos.
Si bien existen diferencias en los tiempos de entrenamiento y
b\'usqueda mayormente, el rendimiento de \texttt{orange-jumpsuit} con conjuntos
de datos peque\~nos se aproxima al de FAISS, lo que destaca la efectividad
de las optimizaciones mediante OpenBLAS y el manejo de memoria alineada.

Finalmente, este trabajo no solo cumple con la implementaci\'on de un sistema
de b\'usqueda eficiente, sino que establece una base s\'olida para futuras
expansiones, como la gesti\'on din\'amica de vectores o el uso de estructuras
de datos m\'as complejas para la inserci\'on ordenada. 

